\chapter{Quantitative Researcher}\label{ch:in:quant-researcher}

A researcher whose signal trades every day and earns a Sharpe ratio of 2 can show, within a year, that the result is not
luck. A researcher whose signal rebalances monthly and earns a Sharpe ratio of one-half needs fifteen and a half years of
live results to show the same. The two carry the same title and may sit in the same building, but one learns from the
market every month and the other perhaps twice in a career. This chapter describes the \omterm{def:in:quant-researcher:def}{quantitative researcher} at each
kind of employer, measures the speed at which the role's work can be judged, and follows what a researcher owns, ships and
comes from; the pay evidence of chapter~14 sits beside it.

\begin{center}\footnotesize
\begin{tabular}{@{}p{3.2cm}p{9.4cm}@{}}
\toprule
\multicolumn{2}{@{}l}{\textbf{Role card: \omterm{def:in:quant-researcher:def}{quantitative researcher}}}\\
\midrule
employers & systematic funds, market makers, platforms, banks, asset managers\\
horizon of decisions & seconds to months\\
P\&L & shares it, through credit for signals\\
reports to & head of research or a portfolio manager\\
occupation codes in filings & 13-2099.01, 15-2041, 15-2031; titles with ``quantitative research'', ``quantitative analyst''\\
where the series teaches it & Book~7, chapters~1 and~15; Book~16, chapter~9\\
\bottomrule
\end{tabular}
\end{center}

\section{Three research jobs: high frequency, medium frequency, asset manager}

\begin{definition}[Quantitative researcher]\label{def:in:quant-researcher:def}
A \emph{quantitative researcher}\index{quantitative researcher} is a person whose work is to find, test and maintain
statistical regularities that a firm can trade or use to manage risk, and to turn them into specifications, models or
code that the firm's trading systems run.
\end{definition}

One market maker describes the role in its own words: \omterm{def:in:quant-researcher:def}{quantitative researchers} ``focus on solving complex problems
through deep, independent research. Their work shapes trading strategies, enhances execution, and improves position
management. Some projects start with a quick idea, while others involve longer investigations and iteration.'' Three
settings give the words different content.
\begin{itemize}
\item \textbf{High frequency} (chapter~2): signals over seconds to minutes, in the order book's own data; the work is
close to the trading system, and a change is judged within days on thousands of independent trades (Book~11).
\item \textbf{Medium frequency} (chapters~4 and~6): signals over days to weeks, in prices, fundamentals and alternative
data; a signal joins a portfolio with others and is judged over months (Books~7 and~8).
\item \textbf{Asset management} (chapter~8): factor and allocation models over months to years, where most of the
evidence comes from history, and live results arrive slowly (Book~8).
\end{itemize}
A bank's quantitative analysts work mainly on pricing and risk models rather than on alpha; chapter~18 covers them.

\section{Horizon, data and feedback speed}

\begin{definition}[Feedback speed]\label{def:in:quant-researcher:feedback}
The \emph{feedback speed}\index{feedback speed} of a research role is the rate at which its work produces evidence that
can tell skill from luck: the calendar time needed for a result of stated size to become statistically distinguishable
from zero.
\end{definition}

\begin{method}[Time to significance]\label{met:in:quant-researcher:time}
Let a strategy's true annual Sharpe ratio be $S$, rebalanced $m$ times a year, with independent returns. The estimated
per-period Sharpe ratio has standard error $\sqrt{(1+s^2/2)/n}$ after $n$ periods, where $s=S/\sqrt m$ (Lo, 2002). The
number of periods needed for the estimate to be $z$ standard errors from zero is
\[
n=\frac{z^2\,(1+s^2/2)}{s^2},
\]
and the calendar time is $n/m$ years. With $z=1.96$, $n$ is close to $3.84\,m/S^2$ periods, and the calendar time close to
$3.84/S^2$ years.
\end{method}

The formula says that calendar time depends on the Sharpe ratio, not on the horizon: a Sharpe ratio of 2 needs about a
year, of 1 about four years, of one-half about fifteen. The horizon enters through the Sharpe ratio a signal can reach.
By the fundamental law of active management (Grinold, 1989; Book~7, chapter~15), a signal with skill $\mathrm{IC}$ per
independent bet and $N$ independent bets a period earns $S\approx \mathrm{IC}\sqrt{N m}$: shorter horizons give more bets
a year and so higher Sharpe ratios, for the same skill per bet.

\begin{example}[The same skill at six horizons]\label{ex:in:quant-researcher:law}
Take an illustrative skill of 0.02 per bet over 50 independent bets each rebalance. The time to significance is then about
194 rebalance periods at every horizon, because $n\approx z^2/(\mathrm{IC}^2N)$ does not depend on $m$: 0.7 days at a
one-minute horizon, 43 days at one hour, 0.77 years at one day, 3.73 years at one week, 16.2 years at one month and 48.5
years at one quarter (\cref{fig:in:quant-researcher:feedback}).
\end{example}

\begin{omfigure}
\begin{tikzpicture}
\begin{axis}[width=10.5cm,height=6cm,xmode=log,ymode=log,xmin=0.001,xmax=100,ymin=0.001,ymax=100,
  xlabel={\small rebalance horizon, trading days (log scale)},ylabel={\small years to significance (log scale)},
  tick label style={font=\footnotesize},grid=major,grid style={gray!20},table/col sep=comma,
  legend style={font=\scriptsize,at={(0.5,-0.3)},anchor=north,legend cell align=left}]
\addplot[omDef,thick,mark=*] table[x=days,y=years]{figdata/industry/17-quant-researcher/feedback.csv};
\addplot[only marks,mark=square*,mark size=2.5pt,omThm] table[x=days,y=years]{figdata/industry/17-quant-researcher/named.csv};
\legend{same skill per bet (IC 0.02; 50 bets a period),Sharpe 2 daily; 1 weekly; 0.5 monthly}
\end{axis}
\end{tikzpicture}

\omcaption{Calendar time before a strategy's Sharpe ratio is distinguishable from zero at 95\%, against its rebalance
horizon (one minute to one quarter). The line holds skill per bet fixed; the squares are the chapter's three named
strategies. Skill and breadth are illustrative. Data: \texttt{in\_researcher.law\_curve} and \texttt{named}.}\label{fig:in:quant-researcher:feedback}
\end{omfigure}

\omterm{def:in:quant-researcher:feedback}{Feedback speed} shapes the job. A high-frequency researcher can try many ideas a year and learn which work; a
medium-frequency researcher must rely more on history, on out-of-sample discipline and on the research controls of Book~7,
because the market will not say in time; an asset-management researcher will rarely see a live result that proves a model
before the model has been used for years. It also shapes careers: the researcher whose results are judged fast can build a
record fast (chapter~28), and pay can follow results more closely (chapter~13).

\section{Ownership of P\&L and of credit}

A researcher rarely owns a book. The profit and loss belongs to a portfolio manager, a desk or the firm, and the
researcher is paid for contributing to it. How the contribution is measured is the firm's credit attribution (Book~16,
chapter~9): signals in a combined portfolio earn together, and dividing their joint profit among their authors is a
choice, not a fact. The three settings differ.
\begin{itemize}
\item At a platform (chapter~5) the researcher works for a portfolio manager whose book is judged on its own result; the
researcher's credit is the manager's decision.
\item At a systematic fund (chapter~4) a signal joins a research portfolio run centrally; credit follows the firm's
attribution method, and the firm owns the signal.
\item At a market maker (chapters~2--3) research improves quoting and hedging that many traders use; credit is shared by
design, and pay follows the desk's or the firm's result.
\end{itemize}
A candidate should ask how credit is measured and who decides it; chapter~13's structures sit on top of the answer.

\section{What a researcher ships}

A researcher's product is not a paper. It is a change that the firm's systems run: a signal with its specification, the
data it needs and how it is computed; the tests and the research log that justify it (Book~7, chapter~1); code that meets
the firm's standards or a specification that developers implement (chapter~19); and a monitoring plan that says how the
firm will know if it stops working. In the market maker's words, researchers ``design and optimize trading algorithms''.
Research reviews decide what ships; most of what a researcher tries does not.

\section{Where researchers come from}

The occupational classification that the filings use for most of the role, 13-2099.01 ``Financial Quantitative
Analysts'', is in O*NET's highest preparation band: ``Most of these occupations require graduate
school. For example, they may require a master's degree, and some require a Ph.D.'' In fiscal 2025, 1\,018 of the
applications classified to the role family carried that code, 156 statisticians (15-2041) and 52 operations research
analysts (15-2031); the role draws on mathematics, statistics, physics, computer science and economics, and chapter~29
follows the pipelines.

\begin{dated}{2025-09}{What quantitative researchers are paid: the public evidence}\label{dat:in:quant-researcher:pay}
US labour condition applications, fiscal 2025, role family ``\omterm{def:in:quant-researcher:def}{quantitative researcher} or analyst'' (titles, or SOC
13-2099.01 when the title is generic): median offered base \$220\,000 at systematic funds (120 applications, 5 employers),
\$190\,000 at platforms (45, 3), \$175\,000 at market makers (131, 9), \$158\,100 at banks (825, 9) and \$107\,100 at
exchanges (35, 4). By \omterm{def:in:pay-levels-by-role-firm-type-and-seniority:soc}{wage level} at the banks: \$88\,300 (I), \$145\,300 (II), \$179\,335 (III), \$200\,000 (IV); at the
market makers \$150\,000, \$200\,000, \$150\,000, \$225\,000. The occupational survey for financial specialists in the
securities industry, May 2025: median \$100\,190, 10th--90th percentile \$55\,370--222\,100. \omterm{def:in:how-pay-works:base}{Base salary} only.
\end{dated}

\begin{omfigure}
\begin{tikzpicture}
\begin{axis}[width=10.5cm,height=5.2cm,xmin=50,xmax=300,ymin=0.4,ymax=5.6,ytick=data,
  yticklabels from table={figdata/industry/17-quant-researcher/pay_kind.csv}{kind},yticklabel style={font=\footnotesize},
  xticklabel style={font=\footnotesize},xlabel={\small offered base, \$ thousand},grid=major,grid style={gray!20},
  table/col sep=comma]
\addplot[draw=none,mark=none,error bars/.cd,x dir=both,x explicit,error mark=none,error bar style={omDef,line width=1pt}]
  table[x=p50,y=pos,x error plus expr=\thisrow{p90}-\thisrow{p50},x error minus expr=\thisrow{p50}-\thisrow{p10}]
  {figdata/industry/17-quant-researcher/pay_kind.csv};
\addplot[draw=none,mark=none,error bars/.cd,x dir=both,x explicit,error mark=none,error bar style={omDef!40,line width=7pt}]
  table[x=p50,y=pos,x error plus expr=\thisrow{p75}-\thisrow{p50},x error minus expr=\thisrow{p50}-\thisrow{p25}]
  {figdata/industry/17-quant-researcher/pay_kind.csv};
\addplot[only marks,mark=|,mark size=5pt,line width=1.5pt,omThm] table[x=p50,y=pos]{figdata/industry/17-quant-researcher/pay_kind.csv};
\end{axis}
\end{tikzpicture}

\omcaption{\omterm{def:in:quant-researcher:def}{Quantitative researchers}' offered base in fiscal 2025 labour condition applications by kind of employer: 10th to
90th percentile (thin), interquartile range (thick), median (mark). Data: \texttt{data/industry/lca\_ranges.csv}, through
\texttt{in\_researcher.pay}.}\label{fig:in:quant-researcher:pay}
\end{omfigure}

The order of employers matches chapter~14's for all roles. The banks' filings show the clearest seniority ladder
(\cref{fig:in:quant-researcher:levels}); the market makers' do not, for the reason chapter~14 gave: offers far above the
\omterm{def:in:pay-levels-by-role-firm-type-and-seniority:lca}{prevailing wage} are not ordered by the level the employer states.

\begin{omfigure}
\begin{tikzpicture}
\begin{axis}[width=10.5cm,height=5cm,xmin=0.8,xmax=4.2,ymin=60,ymax=250,xtick={1,2,3,4},xticklabels={I,II,III,IV},
  xlabel={\small wage level},ylabel={\small median offered base, \$ thousand},tick label style={font=\footnotesize},
  grid=major,grid style={gray!20},table/col sep=comma,unbounded coords=jump,
  legend style={legend cell align=left,font=\scriptsize,at={(0.5,-0.3)},anchor=north,/tikz/every even column/.append style={column sep=8pt}},
  legend columns=3]
\addplot[omThm,thick,mark=square*] table[x=level,y=market_maker]{figdata/industry/17-quant-researcher/levels.csv};
\addplot[omDef,thick,mark=*] table[x=level,y=bank]{figdata/industry/17-quant-researcher/levels.csv};
\addplot[black!70,thick,mark=triangle*] table[x=level,y=systematic_fund]{figdata/industry/17-quant-researcher/levels.csv};
\legend{market maker,bank,systematic fund}
\end{axis}
\end{tikzpicture}

\omcaption{Median offered base of \omterm{def:in:quant-researcher:def}{quantitative researchers} by \omterm{def:in:pay-levels-by-role-firm-type-and-seniority:soc}{wage level}, fiscal 2025. Suppressed cells are left out.
Data: as \cref{fig:in:quant-researcher:pay}.}\label{fig:in:quant-researcher:levels}
\end{omfigure}

\section{Tutorial: how long until you know?}\label{tut:in:quant-researcher:tut}

\begin{tutorial}
\textbf{Goal.} Put numbers on the role's \omterm{def:in:quant-researcher:feedback}{feedback speed} and beside its pay evidence.
\textbf{End state:} \cref{fig:in:quant-researcher:feedback,fig:in:quant-researcher:pay,fig:in:quant-researcher:levels}.
\begin{tutsteps}
\item \textbf{Pay.} \texttt{firm.roles.lca\_cells(rows, role)} with the role \texttt{quant researcher}, at all levels and by
level; \texttt{data/industry/lca\_titles.csv} gives the occupation codes behind the role.
\item \textbf{Time to significance.} \texttt{periods\_to\_significance(S, m)} and \texttt{years\_to\_significance}
(\cref{lst:in:quant-researcher:time}).
\omcode{code/firm/roles/firm_roles.py}{124}{132}{Periods and years of data before a Sharpe ratio is distinguishable from
zero, with Lo's standard error.}\label{lst:in:quant-researcher:time}
\item \textbf{The law.} \texttt{fundamental\_law\_sr(IC, N, m)} turns skill and breadth into a Sharpe ratio for each
horizon.
\end{tutsteps}
\end{tutorial}

For the chapter's three strategies: a Sharpe ratio of 2 at a daily horizon needs 0.97 years (244 trading days); 1 at a
weekly horizon, 3.88 years (202 weeks); one-half at a monthly horizon, 15.53 years (186 months).

\textbf{What to change next.} Let returns be autocorrelated and use Lo's adjusted standard error; count the trials a
researcher ran and deflate the Sharpe ratio for them (Book~7).

\section{Build: the researcher's card and feedback speed}\label{bld:in:quant-researcher:roles}

\begin{build}
\bfield{Purpose} Add the \omterm{def:in:quant-researcher:def}{quantitative researcher} to the role registry, and the functions that measure how fast research is
judged.
\bfield{Interface} \texttt{firm.roles}: the \texttt{quant researcher} card; \texttt{periods\_to\_significance(sr,
periods\_per\_year, z)}; \texttt{years\_to\_significance}; \texttt{fundamental\_law\_sr(ic, breadth,
periods\_per\_year)}.
\bfield{Rules} Returns are IID unless the caller adjusts; the Sharpe ratio is annual; skill and breadth are the caller's,
labelled illustrative.
\bfield{Acceptance tests} \texttt{code/firm/roles/tests/}: the time to significance is $3.84/S^2$ years to first order and
exactly the formula's value; it falls with the Sharpe ratio; under the fundamental law the number of periods is the same at
every horizon.
\bfield{Stretch} Autocorrelation and non-normal returns; multiple-testing deflation; a Bayesian version that updates a
prior on skill.
\end{build}

\begin{omsources}
\item Lo, A.~W. (2002), The statistics of Sharpe ratios, \emph{Financial Analysts Journal} 58(4), 36--52.
\item Grinold, R.~C. (1989), The fundamental law of active management, \emph{Journal of Portfolio Management} 15(3), 30--37.
\item Optiver, ``Breaking down the trading industry''; O*NET, 13-2099.01.
\item Chapter~14's tables from the Department of Labor's LCA files and the BLS occupational survey.
\end{omsources}

\section{Exercises}

\begin{exercise}[$\star$]\label{exo:in:quant-researcher:1}
How many years of daily returns are needed before a true Sharpe ratio of 1.5 is distinguishable from zero at 95\%?
\end{exercise}

\begin{exercise}[$\star$]\label{exo:in:quant-researcher:2}
Under the fundamental law with skill 0.02 and 50 bets a period, what is the annual Sharpe ratio at a weekly horizon?
\end{exercise}

\begin{exercise}[$\star$]\label{exo:in:quant-researcher:3}
From the dated box, which kind of employer offers \omterm{def:in:quant-researcher:def}{quantitative researchers} the highest median base, and by how much does it
exceed the banks'?
\end{exercise}

\begin{exercise}[$\star\star$]\label{exo:in:quant-researcher:4}
Show that under the fundamental law the number of periods to significance does not depend on the horizon.
\end{exercise}

\begin{exercise}[$\star\star$]\label{exo:in:quant-researcher:5}
A researcher's monthly signal has shown a Sharpe ratio of 0.8 over three years of live trading. Is that evidence of skill?
\end{exercise}

\begin{exercise}[$\star\star$]\label{exo:in:quant-researcher:6}
Why might a platform researcher's credit depend more on one person's judgement than a systematic fund researcher's?
\end{exercise}

\begin{exercise}[$\star\star\star$]\label{exo:in:quant-researcher:7}
\emph{Coding.} Returns with first-order autocorrelation $\rho$ inflate the variance of the annualised Sharpe ratio by about
$(1+\rho)/(1-\rho)$. Recompute the time to significance of the weekly strategy with $\rho=0.2$.
\end{exercise}

\begin{exercise}[$\star\star\star$]\label{exo:in:quant-researcher:8}
\emph{Find the flaw.} ``Our monthly strategy has a backtest Sharpe ratio of 1.2 over 20 years, so it will be proven live
within three years.''
\end{exercise}

\section{Problem: How Long until You Know?}

\begin{problem}[{Weekend problem --- how long until you know?}]\label{pb:in:quant-researcher:1}
A graduate with offers in high-frequency research and in a monthly-rebalancing systematic fund asks how quickly each job
would tell her whether she is good at it.

\medskip\noindent\textbf{Part I --- The role.}
\begin{enumerate}
\item Define a \omterm{def:in:quant-researcher:def}{quantitative researcher}; describe the three research settings.
\item How does a market maker describe the role?
\item How does a bank's quantitative work differ?
\item Define \omterm{def:in:quant-researcher:feedback}{feedback speed}.
\item State the role card.
\end{enumerate}

\medskip\noindent\textbf{Part II --- The statistics.}
\begin{enumerate}[resume]
\item State Lo's standard error of the Sharpe ratio under IID returns.
\item Derive the number of periods to significance.
\item Why does calendar time depend on the Sharpe ratio rather than the horizon?
\item State the fundamental law and what it implies about horizon.
\item Give the times for the six horizons of the example.
\end{enumerate}

\medskip\noindent\textbf{Part III --- The job.}
\begin{enumerate}[resume]
\item Who owns the P\&L at a platform, a systematic fund and a market maker?
\item What does a researcher ship?
\item Where do researchers come from, and what preparation does the occupation require?
\item What does the pay evidence show by employer and by level?
\item Why do the market makers' filings not show a seniority ladder?
\end{enumerate}

\medskip\noindent\textbf{Part IV --- The verdict.}
\begin{enumerate}[resume]
\item State the \emph{named result}: the time to reject a zero Sharpe ratio for a Sharpe ratio of 2 at a one-day horizon,
1 at one week and 0.5 at one month.
\item What would autocorrelation do to these times?
\item What can a researcher in a slow-feedback job do to learn faster?
\item How should the difference affect how she judges her first two years in each job?
\item In two sentences, answer her.
\end{enumerate}
\end{problem}

\section{Interview questions}

\begin{interviewq}[$\star$ \iqroles{researcher}]\label{iq:in:quant-researcher:1}
A strategy's daily Sharpe ratio estimate is 0.1 over 400 days. What is its annualised value, and roughly its standard error?
\end{interviewq}

\begin{interviewq}[$\star$ \iqroles{researcher, trader}]\label{iq:in:quant-researcher:2}
Why do high-frequency strategies often have high Sharpe ratios and low capacity?
\end{interviewq}

\begin{interviewq}[$\star\star$ \iqroles{researcher}]\label{iq:in:quant-researcher:3}
You tried 200 signals and the best has a backtest t-statistic of 3.5. What do you conclude?
\end{interviewq}

\begin{interviewq}[$\star\star$ \iqroles{researcher, mle}]\label{iq:in:quant-researcher:4}
How would you decide that a live signal has stopped working?
\end{interviewq}

\begin{interviewq}[$\star\star$ \iqroles{researcher, developer}]\label{iq:in:quant-researcher:5}
What must a signal's specification contain so that a developer can implement it without you?
\end{interviewq}

\begin{interviewq}[$\star\star\star$ \iqroles{researcher}]\label{iq:in:quant-researcher:6}
Two signals share 60\% of their positions. How would you divide the credit for their joint profit?
\end{interviewq}
